% ===== §3 Ontology: relation, grammar, value =====
\section{The Ontology of the Generative System}
\label{sec:ontology}

\noindent
The phenomenology determines the task of the ontology. The task is not the general specification of what value or relation is, but the specification of what they must be in order that the datum of \xref{sec:phenomenology} be possible, a datum in which a nameless value is co-generated in co-apprehension, conditioned by the manner of a joint attending directed upon an opening. The ontology's theses are warranted as the conditions of that datum, and are not advanced beyond it. The section establishes three theses concerning the being of the generative system, each as a condition of the founding datum rather than as an independent posit, and concludes with the single structural feature, reflexivity, of which the two following sections supply distinct disciplinary realisations.

\subsection{Relation, grammar, and value as aspects of one system}
\label{sec:three-aspects}

The construal of relation, grammar, and value as three entities standing in relations of causal or constitutive dependence, with a consequent question as to which is fundamental, is to be declined. That construal imports the externality the datum excludes: it posits an order of grounding among three separate items and thereby requires a standpoint external to all three from which their order is surveyed. The phenomenology affords no such standpoint. It affords a single event under three aspects, and the ontology specifies the aspects without dividing the event.

\begin{claim}[Relation, grammar, and value as three aspects of one generative system]
Relation, grammar, and value are not three entities in relations of mutual foundation but three ontological aspects of a single generative system. Relation is the being of the system: the relational substance, that which is, not a phenomenon of a deeper term and not a representation. Grammar is the manifestation of the system's dynamics: not a description applied externally to the system in order to cognise it, but the form in which the system's dynamics manifests itself, the morphology of its movement. Value is the phenomenon of the system: the aspect of the generative system that emerges and admits of being witnessed. As aspects of one event rather than items in a sequence, the three are co-temporal of necessity; there is no moment at which the relation obtains while the grammar and value do not, since they are one being under three regards.
\end{claim}

\noindent
This dissolves the circularity that would otherwise obtain. The objection that relation, value, and grammar each presuppose the others, so that none can be initiated, assumes them to be distinct items requiring sequence. They are not distinct items. They are the being (relation), the manifest dynamics (grammar), and the appearance (value) of one generative system, and one entity under three regards requires no sequence, as a body's being, its motion, and its visible aspect are not three items awaiting one another but one entity regarded in three ways. The three registers of value established in the prior paper, the imaginary, the symbolic, the real, are accordingly relocated as distinctions internal to the third aspect, articulations within the aspect termed phenomenon, and not as competitors to relation or grammar.

\subsection{The ontological status of grammar}
\label{sec:grammar-ontological}

A consequence is to be drawn, as it governs the subsequent development. Grammar, in this ontology, is an ontological category and not an epistemic one. It is not the formal apparatus constructed by an observer in order to describe or predict the relation; that apparatus, the externalist's grammar, is a settling, a fixing of the dynamics into a surveyable rule-set, the operation \xref{sec:grammar} identifies as a form of possession. The grammar at issue is the self-manifestation of the dynamics: the form the generative movement assumes in moving, antecedent to and independent of any observer's attempt to fix it. The characteristic rhythm of two parties' being-together, the moves that are available to them and those that are not, the developing repertoire of private signs, is not an observer's description of their relation but the self-manifestation of the relation's dynamics. The grammar is in the movement, not in the notation.

This distinction supports a subsequent result. It permits the consistent assertion both that the relational grammar is real, being the system's own manifest dynamics, and that no notated grammar captures it, every notated grammar being an external settling that, in fixing the dynamics, falsifies the feature, open and continuing generation, constitutive of it. The grammar is, ontologically; it is non-possessible, epistemically; and \xref{sec:grammar} establishes that this non-possessibility is not a limitation of instruments but the formal signature of the relation's non-possessibility.

\subsection{The reflexive structure}
\label{sec:reflexive}

The structural feature that organises the paper, located phenomenologically in \xref{sec:hinge} and realised in distinct forms in the two following sections, is now stated ontologically. In a standard dynamical system the law of evolution is fixed: the state evolves, while the rule governing its evolution does not change in consequence of the evolution. The datum of \xref{sec:phenomenology} departs from this. There, apprehension was generation: the coming-to-be of the value was not the output of a fixed dynamics but an event whose manner conditioned what could subsequently obtain. Stated ontologically:

\begin{claim}[The reflexive reconfiguration of the dynamics by the phenomenon]
The value, which is the phenomenon, does not merely issue from the grammar, which is the manifest dynamics; it reconfigures it. Each crystallisation of a nameless value into the relation alters the system's dynamics, modifying which moves are subsequently available, opening regions of the relational space that did not antecedently exist and foreclosing others. The generative system is one whose dynamics is reconfigured by its own products: the appearance reacts upon and reconfigures the movement that bore it. The system is therefore not one with a fixed law and a varying state, but one in which the law is itself transformed by what it produces, a self-modifying generativity.
\end{claim}

\noindent
Three consequences follow, each of which is discharged in the subsequent development.

The \emph{first} is that the subject is generated within the cycle and is not antecedent to it. If the dynamics is reconfigured by its products, and if the parties to the relation are among the system's products, constituted, as the prior treatment held, in and through the relation rather than entering it complete, then the subject is not an external term standing outside the generative cycle and operating it. The subject is a product of the cycle it contributes to generating: a fixed point of the reflexive structure, not its external premise. There is no antecedently complete self that enters the relation in order to generate value; the self is, in part, what the relation's value-cycle generates. This is the antecedent that \xref{sec:political} develops into a political economy, the thesis that the circulation of value reconfigures the subject who creates it being Marx's, transposed here from the social to the ontological register.

The \emph{second} is that the dynamics admits no external pre-enumeration. A system whose law is reconfigured by its products has no fixed law that an external observer could specify in advance of the system's operation, the subsequent law being generated within the operation, in the joint act, and not antecedently available. This is the ontological antecedent of the non-enumerability that \xref{sec:grammar} formalises: not the thesis that a super-physical process operates, but the thesis that any external fixing of the grammar fails because the grammar fixed is in the course of being generated, internally, by the generation the fixing would be required to anticipate.

The \emph{third} is that the manner conditions the product without remainder. Since apprehension is generation and the dynamics is reconfigured by what it generates, the manner of the system's operation, the void around which, the register within which, the quality with which the opening is sustained, is not incidental to the product. There is no separation of a neutral content of value from the manner of its co-generation; the manner is constitutive. This is the antecedent of the result that the ethics of \xref{sec:keystone} is not an ethics of externally assessed outcomes, and that the practice of \xref{sec:praxis} is the tending of the manner, of the field and its curvature, rather than the manufacture of the product.

\subsection{The commitments and exclusions of the ontology}
\label{sec:onto-refuses}

The character of the ontology is to be stated, together with its exclusions. It is a process ontology, and a dialectical one: there is no value antecedent to the relation's movement, and no movement antecedent to the relation; there is one generative system, manifest as relation, grammar, and value, whose dynamics is transformed by its own appearing. It is continuous in this respect with the relational ontology of the mother paper, being is relational, the terms are constituted in the relating, and adds the reflexive clause, that the relating is reconfigured by what it generates.

The ontology excludes three theses, and the exclusions are of equal standing with the commitments. It excludes substance: there is no value-substance, no relational substrate, subsisting beneath the system and persisting unchanged through its transformations; the three aspects are aspects of a movement, not properties of a thing. It excludes the external standpoint: there is no position external to the generative system from which its three aspects are surveyed and their order legislated, which is the ontological form of the absence of a metalanguage that the prior paper established, and which \xref{sec:breaking} discharges in addressing the cognition of the unbroken structure. And it excludes completion: the system is not a totality that is, but a totality in generation, its law in the course of being written; to treat it as complete, to notate its grammar and declare it captured, is the ontological form of settling, and, as \xref{sec:keystone} establishes, of alienation. These exclusions are not concessions. They are the conditions of the founding datum: a value antecedent as substance could not be co-generated; a system surveyable externally would not require the joint internal witnessing of which the datum consists; and a completed totality could generate nothing, least of all a value not yet named.

The ontology possesses a single structural signature, reflexivity: the reconfiguration of the dynamics by the phenomenon, of the producer by the product. The signature is, to this point, abstract, stated in the general terms of relation, grammar, and value. The two following sections supply two disciplinary realisations of this structure, the abstract structure being realised, under the conditions of distinct disciplines, in distinct and mutually irreducible forms. Under the conditions of formal epistemology it is realised as a participatory, quantum-structural grammar (\xref{sec:grammar}); under the conditions of political economy it is realised as the circulation of value that reconfigures its subject (\xref{sec:political}). Neither realisation is the truth of the other, and neither is the unbroken structure; they are two symmetry-broken phases of it, and the cognition of the unbroken structure they share, in the absence of the external standpoint the ontology excludes, is the matter of \xref{sec:breaking}.
